\section{Introduction}\label{sec:current-overview}\label{sec:introduction}
A predictive state records everything about a past experiment that can
affect its future laws. Its linear dimension does not, by itself, determine
the memory needed to realize those laws with nonnegative probabilities.
Commands can change the next observation, an observation changes the state
on which later commands act, and retaining or communicating information
uses resources. This paper studies the resulting positive-realization
problem within a theory of causal resource comparison and optimal
statistical discrimination.

We distinguish two response objects. A terminal response assigns a
probability row $b(w)$ to each finite command word. An instrument response
assigns a joint probability to each finite command--outcome history. An
exact classical realization uses one fixed collection of stochastic tables
on finitely many live labels. Except in an explicitly stated mixed-start
comparison, the initial label is pure and counted. All lower bounds allow
arbitrary real probabilities. Program-table length, arithmetic precision,
working storage and the protected physical system are separate resources.

The central questions are structural. Which linear predictive actions
admit finite positive realizations? When does linear representation theory
determine the exact resource, and when do probability constraints add a
strict cost? How does the preparation contribute to that cost? Can one
controller reproduce a whole commanded instrument, rather than a scalar
subsequence of its observations? The results below answer these questions
on distinct, explicitly specified classes.

\subsection*{Reversible responses and finite-group structure}
Let $V_b$ be the real span of residual functions $b^w(z)=b(wz)$.
If command shifts are invertible on $V_b$ and generate a group with compact
closure $G_b$, Theorem~\ref{thm:r72-reversible} characterizes finite exact
classical realizability by finiteness of $G_b$. This applies to prepared
finite-dimensional unitary commands after invisible directions have been
removed. Theorem~\ref{thm:r73-normal-form} gives more than a finiteness
criterion: every finite stochastic realization reduces, without increasing
its size, to a permutation core and at most one counted initializer.
The original hidden transition matrices need not form a group. If $N$ is
the number of distinct residuals and $\kappa_G(b)$ is the least size of a
probability factorization $b(g)=\lambda P_gD$, then
\begin{equation}\label{eq:r75-intro-core}
 m(b)=\min\{N,1+\kappa_G(b)\}.
\end{equation}
Here both the preparation $\lambda$ and every decoder row are probability
vectors. Their positivity constraints cannot be replaced by a rank test.

Let $\pi$ be the group average and put $b_t=\pi+t(b-\pi)$.
The retained Theorem~\ref{thm:r74-low-contrast} proves, for every finite
observable group and finite output alphabet, that sufficiently small
positive contrast makes $\kappa_G(b_t)$ equal to the smallest permutation
module containing $V_b$. The new Theorem~\ref{thm:r75-orbit-multiplicity}
turns that minimum into the finite integer problem
\[
 \nu_G(V_b)=\min\left\{\sum_{[H]}q_H[G:H]:
   \sum_{[H]}q_H\dim(\rho^H)\ge r_\rho(b)
   \text{ for every }\rho\in\widehat G\right\}.
\]
The variables are nonnegative integers, $[H]$ ranges over conjugacy classes
of subgroups, and $r_\rho(b)$ is the multiplicity of the complex irreducible
$\rho$ in the complexified response module. Stacked Fourier blocks compute
these multiplicities. The proof includes descent to a real equivariant
embedding, so no real or quaternionic representation type is omitted.
For a prescribed transitive action $H\backslash G$, the same argument gives
the exact number of identical orbit copies needed at small positive
contrast. These are exact resource formulas, without a polynomial-time
claim for subgroup enumeration or integer optimization.

The probability constraint also singles out a precise algebraic boundary.
Theorem~\ref{thm:r75-normality-convexity} proves that the set of binary
responses realizable on one transitive orbit $H\backslash G$ is convex if
and only if $H$ is normal. The forward obstruction is explicit: averaging
the deterministic return responses gives the normalized permutation
character. If $H$ is not normal, one Fourier block has larger rank than
any single-orbit response can have. This obstruction persists at every
positive contrast toward the constant average. When $H$ is normal,
preparation can be absorbed into a quotient decoder, as in the retained
Theorem~\ref{thm:r74-quotient-lp}. Thus the earlier quotient linear programs
have an exact structural explanation. The result concerns this orbit
response set; it does not exclude other linear formulations of special
nonnormal problems.

Probability constraints yield a further, quantitative rigidity. The sharp
singular-matrix bound in Lemma~\ref{lem:r75-singular-coefficient}, combined
with subgroup-index estimates, gives a complete contrast law for an
unbounded nonabelian family. For $G=S_n$, $n\ge5$, let
$\chi_n(g)$ count fixed points of the natural action and define
\[
 \beta_t(g)=\frac{1-t}{n}+\frac{t}{n}\chi_n(g),\qquad
 b_t(g)=(\beta_t(g),1-\beta_t(g)).
\]
Put $\tau_n=(\sqrt n-1)/(\sqrt n+1)$. Theorem~\ref{thm:r75-symmetric-return}
proves, over all finite stochastic controllers,
\[
 m(b_t)=
 \begin{cases}
 1,&t=0,\\
 n(n-1)+1,&0<t\le\tau_n,\\
 n^2+1,&\tau_n<t\le1.
 \end{cases}
\]
Throughout positive contrast, the group is $S_n$, the residual count is
$n!$, and the predictive dimension is $1+(n-1)^2$. The lower bound covers
every finite stochastic controller through \eqref{eq:r75-intro-core}, not
only a chosen collection of natural orbits. The upper tables are explicit;
two command types and binary readout suffice for every $n$. At the
threshold, the coefficient matrix is unique up to coordinate permutation.
Its minimum permutation-core factorizations have fixed preparation rows
and an explicitly described freedom in their orbit weights and decoder
rescaling. This classifies these optimal cores, without asserting that
all hidden stochastic realizations are unique. At full contrast, the
identity value and group mean give the separate point-return rigidity of
Lemma~\ref{lem:r75-return-rigidity}. A diagonal quantum realization supplies
a physical implementation without a quantum-advantage assertion.
The earlier sharp $S_3$ and four-bit examples and the quotient payoff
duals are retained as separate complete model laws.

\subsection*{Initialization, conditional dynamics and one common simulator}
Preparation has a general cost beyond reversible experiments. For any
stationary process over a finite alphabet, a pure hidden initial state has
the same future-response function as a convex combination of its
one-step successors. Unless the process is independent and identically
distributed, that initial state can be eliminated from any pure-start
realization. The resulting exact identity is
\[
 n_{\mathrm{pure}}=n_{\mathrm{mixed}}+1
\]
for every non-i.i.d. stationary process admitting a finite classical model.
Theorem~\ref{thm:r75-stationary-initializer} proves this by an explicit state elimination, including
preservation of every word law. The i.i.d. case has both minima equal to
one; a process with no finite mixed realization has no finite pure one.
A controlled version applies when the predictive preparation is invariant
under a full-support average of unobserved commands. This explains the
initialization charge by stationarity itself, independently of a special
phase or a minimal-rank assumption.

The positive physical mechanism remains the retained
Theorem~\ref{thm:r68-instrument}. For arbitrary known finite instrument
families on $\mathbb C^h$, a uniform lower bound $\eta I/h$ on every
normalized branch yields a common finite physical dictionary. Branchwise
barycentric expansions give one set of nonnegative tables for all commands,
all visible adaptive histories, terminal POVMs, complete stopped laws and
nonnegative history costs. The theorem also provides sparse branch rows,
several-step variants and its stated state-positive noise exponent.
These are inherited results, with their complete proofs in the supplement.

Average mixing alone has a different effect. The retained
Theorem~\ref{thm:r74-depolarizing-obstruction} gives a stationary,
one-dependent qubit process with Hankel rank three, uniformly positive
outcome probabilities and an exactly depolarizing unconditioned channel.
Its exact classical memory is $q$ with a mixed start and $q+1$ with a counted
pure start for a conditional phase of order $q\ge3$, and is infinite for a
phase of infinite order. Theorem~\ref{thm:r75-noisy-memory} develops the corresponding
noisy \emph{complete instrument}. One regular polygon is simultaneously
invariant under all three observed branches. The reset branch still
prepares a pure state at every noise value, so the construction does not
follow by imposing a positive interior margin on all posteriors, even
after a fixed block of calls. The relative spectral lower
bound and this common construction give exact first-sector memory criteria
and the exact worst-phase count
\[
 \max_\theta n_{\mathrm{mixed}}=K(\alpha)
   :=\max\{3,\lceil\pi/\arccos\alpha\rceil\},\qquad 0<\alpha<1.
\]
The scalar sector and polygon geometry were established in the earlier
spectral theory. The additional conclusion here is that its bounds are
attained for the full observed instrument, with one physical dictionary,
and that initialization adds exactly one label by the general theorem.
The zero-retention endpoint $\alpha=0$ has two mixed and three pure
states and is treated separately.

For $0<\alpha<1$ and any integer $K\ge K(\alpha)$, one $K$-vertex
dictionary works for every rotation command in any supplied finite
alphabet. Its tables reproduce all adaptive command--outcome laws. For
$K=K(\alpha)$, the bound is exact when the alphabet includes a maximizing
phase from Theorem~\ref{thm:r75-noisy-memory}. This quantifier differs from choosing a new simulator
for each scalar phase. Theorem~\ref{thm:r75-stopped-approximation} perturbs only the rotating branch and compensates
the third branch to give one finite simulator for the undamped experiment
whose stopped-record error is at most $(1-\alpha)\mathbb E_1[\tau]/8$.
For policies whose undamped target expected calls are at most $B$,
and error tolerance $\delta$, this gives an
explicit $O(\sqrt{B/\delta})$ label upper bound, uniformly over visible
adaptive stopping and command policies. The approximate experiment need
not retain the target expected-duration bound. It does not establish a matching
approximation lower bound or solve a different minimax-risk frontier.

\subsection*{Priority, dependence and the original causal theory}
Residual automata, stochastic realization, Hankel rank, invariant cones and
observable operator models are classical
\cite{RabinScott59,Heller65,CarlylePaz71,Jaeger00,ThonJaeger15,MonrasWinter14}.
The structure of stochastic matrix groups
\cite{Schwarz64,Flor69,Brown64,Schwarz62}, finite-group representation theory
\cite{Serre77}, symmetric positive factorizations
\cite{GPT13,KPT10} and nonnegative-factorization perturbations
\cite{BocciCarliniRapallo11,GillisGlineur12} supply established ingredients.
We credit multiplicity decomposition, Frobenius reciprocity and Fourier
orthogonality as classical facts, with the steps needed here proved below.
The new group assertions are the exact orbit-convexity boundary, its
positive resource interpretation, and the sharp symmetric-group rigidity
law. They extend the already proved reversible normal form and contrast
construction rather than replacing them.

Invariant-polytope simulation of quantum instruments, including adaptive
composition, has explicit antecedents \cite{MonrasWinter14,Okay26}.
Fanizza, Lumbreras and Winter \cite{Fanizza24} treat phase obstructions,
finite quantum processes without finite classical realizations and noisy
classical-memory lower bounds. One-dependent non-HMM processes also
predate this work \cite{HolroydLiggett16}. Neither phenomenon is claimed
as a first. The present assertions concern the general counted-initialization
identity and the displayed qubit family's exact, simultaneous full-instrument
resource laws and stopped-law approximation. Finite diagnostics below are
checks of identities and constructions, not substitutes for the proofs or
for a complete literature-priority assessment.

The main article retains the full actual-history comparison theorem and
its common paid construction. Its witness is evaluated on the competitor's
actual pre-call histories. The paired supplement retains performed-law
compactness, variational and dual attainment, robust tied optimality,
normalized stopping rigidity, finite-width transfer, rank-constrained
optimal discrimination, finite-order quantum measurement, two-stage reset
control, memory and communication accounts, spectral geometry and strategy
implementation. The independently compiled companion retains absolute
calibration and nonintervening directional sensing.

Each layer keeps its own quantifiers: general causal comparison for its
declared experiment class, exact reversible realization under compact
reversibility, stationary initialization for arbitrary stationary processes,
interior-posterior simulation for its stated instrument class, and closed
model laws for their specified data. Local geometric conclusions remain
local. None completes the separate A1, A2 or further Theta construction
objectives. The current theorem and preservation maps distinguish logical
proof dependencies from retained historical material.
